\section{Mathematical Frameworks for LLM Evaluation}
\label{sec:method}

We propose five evaluation frameworks, each grounded in established mathematical theory and addressing a specific gap in current methodology.
We first describe the common experimental setup, then present each framework with formal definitions and algorithmic descriptions.

\subsection{Experimental Setup}
\label{sec:setup}

\para{Data.}
We use the \mmlu dataset~\citep{hendrycks2021mmlu} comprising 14,042 test questions across 57 subjects spanning STEM, humanities, social sciences, and professional domains.
We construct performance profiles for 12 models based on published accuracy scores: \gpt~\citep{openai2023gpt4}, GPT-4o~\citep{openai2024o1}, \claude Opus and Claude-3.5 Sonnet~\citep{anthropic2024claude3}, \llama 8B and 70B~\citep{dubey2024llama3}, Mistral 7B, Qwen-2 72B, \gemini 1.5 Pro~\citep{geminiteam2024gemini15}, PaLM-2, \mixtral~\citep{jiang2024mixtral}, and Phi-3 Medium.

\para{Representation.}
Each model $m_i$ is represented as a capability vector $\vc_i \in [0,1]^{57}$, where coordinate $j$ denotes accuracy on subject $j$.
The collection $\gC = \{\vc_1, \ldots, \vc_{12}\} \subset \sR^{57}$ forms a point cloud amenable to topological and geometric analysis.

\para{Implementation.}
All frameworks are implemented in Python using NumPy, SciPy, Ripser~\citep{bauer2021ripser} for persistent homology, Persim for persistence diagram distances, NetworkX for graph construction, and Scikit-learn for embeddings.

\subsection{Framework 1: Topological Drift Detection via Persistent Homology}
\label{sec:topdrift}

\para{Motivation.}
Current evaluation metrics provide point estimates without formal guarantees about how rankings change under perturbation.
We apply \ph to detect structural changes in model capability landscapes across time.

\begin{definition}[Capability Point Cloud]
Given models $\{m_1, \ldots, m_n\}$ evaluated on $d$ tasks, the \emph{capability point cloud} is $\gC = \{\vc_i\}_{i=1}^n \subset \sR^d$ where $\vc_i = (\text{acc}_{i,1}, \ldots, \text{acc}_{i,d})$.
\end{definition}

\begin{definition}[Vietoris-Rips Filtration]
For parameter $\epsilon \geq 0$, the Vietoris-Rips complex $\text{VR}_\epsilon(\gC)$ contains a simplex $[\vc_{i_0}, \ldots, \vc_{i_k}]$ whenever $\|\vc_{i_p} - \vc_{i_q}\| \leq \epsilon$ for all pairs $p, q$.
The nested family $\{\text{VR}_\epsilon(\gC)\}_{\epsilon \geq 0}$ is the Vietoris-Rips filtration.
\end{definition}

\PH computes the birth-death pairs of topological features (connected components $H_0$, loops $H_1$, cavities $H_2$, etc.) across the filtration, producing a persistence diagram $\text{Dgm}(\gC)$.

\begin{theorem}[Stability~\citep{cohen2007stability}]
\label{thm:stability}
For point clouds $\gC, \gC'$, the bottleneck distance between persistence diagrams satisfies:
\begin{equation}
d_B(\text{Dgm}(\gC), \text{Dgm}(\gC')) \leq d_H(\gC, \gC')
\end{equation}
where $d_H$ denotes the Hausdorff distance.
\end{theorem}

\para{Drift detection.}
We partition models into temporal cohorts: $\gC_{\text{early}}$ (2023 models) and $\gC_{\text{late}}$ (2024 models).
Structural drift is quantified by the Wasserstein-1 distance $W_1(\text{Dgm}(\gC_{\text{early}}), \text{Dgm}(\gC_{\text{late}}))$ between persistence diagrams.
By Theorem~\ref{thm:stability}, this provides a formal lower bound on the minimum geometric change between cohorts:
\begin{equation}
d_H(\gC_{\text{early}}, \gC_{\text{late}}) \geq d_B(\text{Dgm}(\gC_{\text{early}}), \text{Dgm}(\gC_{\text{late}}))
\end{equation}

\subsection{Framework 2: Sheaf-Theoretic Compositional Evaluation}
\label{sec:sheafeval}

\para{Motivation.}
Multi-metric evaluations (\eg \helm) aggregate local scores into global rankings without formal guarantees about information preservation.
We formalize evaluation composition as a sheaf on a task-decomposition graph.

\begin{definition}[Evaluation Sheaf]
Let $G = (V, E)$ be a directed graph where vertices represent evaluation granularities (global, category, subject) and edges represent containment.
A \emph{cellular evaluation sheaf} $\gF$ on $G$ assigns:
\begin{itemize}[leftmargin=*,itemsep=0pt,topsep=0pt]
    \item To each vertex $v \in V$: a vector space $\gF(v) \cong \sR^{n_v}$ of evaluation scores;
    \item To each edge $e: u \to v$: a linear restriction map $\gF_{e}: \gF(u) \to \gF(v)$.
\end{itemize}
\end{definition}

\begin{definition}[Consistency and Cohomology]
A \emph{global section} $s$ assigns to each vertex $v$ a vector $s(v) \in \gF(v)$ such that for all edges $e: u \to v$, $\gF_e(s(u)) = s(v)$.
The \emph{first sheaf cohomology} $H^1(G; \gF)$ measures the obstruction to extending local sections to global ones.
When $\dim H^1 > 0$, local evaluations cannot be consistently aggregated.
\end{definition}

For \mmlu, we construct a three-level graph: one global node, 5 category nodes (STEM, Humanities, Social Sciences, Professional, Other), and 57 subject nodes, yielding $|V| = 63$ nodes and $|E| = 65$ edges.
The restriction maps implement weighted averaging from subjects to categories.

\para{Consistency ratio.}
For each model, we compute the ratio between the residual under restriction (how much information is lost by aggregation) and the local evaluation norm:
\begin{equation}
\rho(m) = \frac{\sum_{e: u \to v} \|\gF_e(s_m(u)) - s_m(v)\|^2}{\sum_{v} \|s_m(v)\|^2}
\end{equation}
Models with high $\rho$ have heterogeneous capability profiles where aggregation is unfaithful.

\subsection{Framework 3: Information-Geometric Evaluation Manifolds}
\label{sec:fisheval}

\para{Motivation.}
Comparing models by accuracy ignores distributional structure.
Two models with identical average accuracy may have fundamentally different capability distributions.
We equip the space of model capability distributions with a Riemannian metric.

\begin{definition}[Capability Distribution]
For model $m$ with capability vector $\vc \in [0,1]^d$, we define the normalized capability distribution $p_m = \vc / \|\vc\|_1$ on the $(d-1)$-simplex $\Delta^{d-1}$.
\end{definition}

\begin{definition}[Fisher-Rao Distance]
The \fr distance between models $m_i$ and $m_j$ on the probability simplex is:
\begin{equation}
d_{FR}(p_i, p_j) = 2 \arccos\left(\sum_{k=1}^d \sqrt{p_i^{(k)} \cdot p_j^{(k)}}\right)
\end{equation}
This is the geodesic distance under the Fisher information metric on the multinomial family.
\end{definition}

The \fr distance satisfies the axioms of a proper metric and is the unique Riemannian metric (up to scale) that is invariant under sufficient statistics~\citep{cencov1982statistical}.
This provides principled model comparison beyond raw accuracy.

\para{Fisher information trace.}
The trace of the Fisher information matrix $\text{tr}(\mI_m) = \sum_k 1/p_m^{(k)}$ measures the total sensitivity of the distribution to perturbation.
Models with concentrated capability (high accuracy on few subjects) exhibit higher Fisher information and are more sensitive to evaluation noise.

\subsection{Framework 4: Failure Mode Landscape via TDA}
\label{sec:failtda}

\para{Motivation.}
Standard evaluation reports aggregate error rates but ignores the \emph{structure} of failures.
Two models with identical error rates may fail on completely different questions for different reasons.
We apply TDA to the space of failure patterns to detect structured failure modes.

\begin{definition}[Failure Correlation Graph]
For threshold $\tau \in [0,1]$, define the failure indicator $f_{i,j} = \mathbb{1}[\text{acc}_{i,j} < \tau]$ for model $i$ on subject $j$.
The failure correlation matrix $\mR \in \sR^{d \times d}$ has entries $R_{jk} = \text{corr}(f_{\cdot,j}, f_{\cdot,k})$ measuring co-failure across models.
\end{definition}

We construct a filtered simplicial complex by thresholding the correlation matrix at increasing values and compute persistent homology of the resulting Vietoris-Rips complex.
The emergence of $H_1$ features (loops) indicates cyclic failure relationships---subjects that form failure chains not decomposable into hierarchical clusters.

\para{Phase transitions.}
We track Betti numbers $\beta_0$ (components) and $\beta_1$ (loops) across thresholds $\tau \in \{0.5, 0.6, 0.7, 0.8, 0.9\}$.
A sharp increase in $\beta_1$ at threshold $\tau^*$ indicates a phase transition where structurally distinct failure modes emerge.

\subsection{Framework 5: Spectral Contamination Detection via Random Matrix Theory}
\label{sec:speccontam}

\para{Motivation.}
Benchmark contamination inflates scores in ways that surface-level methods (n-gram overlap, embedding similarity) cannot detect when data is rephrased~\citep{yang2023rethinking}.
We propose detecting contamination through spectral anomalies in performance matrices.

\begin{definition}[Performance Matrix]
The performance matrix $\mP \in \sR^{n \times d}$ has entry $P_{ij}$ equal to the accuracy of model $i$ on subject $j$.
Under the null hypothesis of no contamination, the sample covariance $\mS = \frac{1}{n}\mP^\top \mP$ should have eigenvalues distributed according to the \marchpast law.
\end{definition}

\begin{theorem}[Marchenko-Pastur Law~\citep{marchenkopastur1967}]
For $\mP$ with i.i.d.\ entries of variance $\sigma^2$ and aspect ratio $\gamma = d/n$, the empirical spectral distribution converges to:
\begin{equation}
\rho_{MP}(\lambda) = \frac{1}{2\pi\sigma^2\gamma}\frac{\sqrt{(\lambda_+ - \lambda)(\lambda - \lambda_-)}}{\lambda}
\end{equation}
where $\lambda_\pm = \sigma^2(1 \pm \sqrt{\gamma})^2$ are the bulk edges.
\end{theorem}

\para{Contamination detection.}
Eigenvalues exceeding $\lambda_+$ indicate structure beyond random variation.
Under contamination, subjects with leaked data exhibit correlated performance boosts, creating additional outlier eigenvalues.
We detect contamination via:
\begin{enumerate}[leftmargin=*,itemsep=0pt,topsep=0pt]
    \item Compute eigenvalues of the sample covariance of the performance matrix.
    \item Identify outliers exceeding \marchpast bounds $[\lambda_-, \lambda_+]$.
    \item Apply the \tw test~\citep{tracy1996orthogonal} to the largest eigenvalue for significance.
    \item Localize contaminated subjects via leverage scores of outlier eigenvectors.
\end{enumerate}

\begin{proposition}[Contamination Signal]
If $k$ subjects are contaminated with correlated performance boost $\delta$, the performance matrix gains an additional rank-$k$ component, producing $k$ eigenvalues exceeding $\lambda_+$ by $\Omega(n\delta^2)$.
\end{proposition}
